\section{Reproduction, package contents, and integration}
\label{sec:reproduce}

The package is self-contained for its default exact tests and
coverage replay.  It contains the full article source, rendered PDF,
vector figures, the new runtime modules, a pinned snapshot of their
maintained dependencies, focused tests, research drivers, a frozen
complete oracle, raw experimental evidence, and an additive patch
against the pinned repository commit.

\begin{center}
\small
\begin{tabularx}{\linewidth}{lX}
\toprule
Path & Contents and purpose\\
\midrule
\code{article/} & Main TeX, included sections, bibliography,
figures, tables, PDF, and figure-generation script.\\
\code{code/} & Standalone native package snapshot, research
drivers, selected fixtures, and the 39-test reproduction set.\\
\code{fixtures/} & Unchanged frozen complete discovery corpus
from the preceding minimum-envelope report.\\
\code{reference/} & Unmodified optimized rank-two envelope
module used as the stronger performance control.\\
\code{evidence/} & Full corpus results, compact summaries,
coverage certificates, replay outputs, family component
certificates, raw timings, and regression logs.\\
\code{scripts/} & Independent genuine/abstract differential
audit and sparse-constructor benchmark.\\
\code{integration.patch} & Additive repository changes;
unchanged dependency snapshots are not part of the patch.\\
\code{MANIFEST.json} & File sizes, SHA-256 digests, and
origin classification for every packaged file except itself.\\
\code{reproduce.py} & Manifest verification and isolated
reproduction commands.\\
\bottomrule
\end{tabularx}
\end{center}

From the package root, run
\begin{lstlisting}
python reproduce.py
\end{lstlisting}
This checks the manifest, runs the 31 new focused tests and eight
maintained sector-integration tests, and independently replays the
thirteen corpus and four family coverage certificates.  It writes
fresh output into a new reproduction directory.  It does not
overwrite the supplied evidence.

The longer computations are explicit options:
\begin{lstlisting}
python reproduce.py --audit --family --differential
python reproduce.py --fresh-regina
python reproduce.py --benchmark --sparse-benchmark
python reproduce.py --pdf
\end{lstlisting}
The first command repeats the full frozen-oracle audit, native
family census, and additional differential checks.  The second
also regenerates the designated complete Regina lists; it needs
Regina 7.4 or a compatible installation.  The third repeats the
paired timing protocol, including resource-limited incumbent runs.
The fourth compiles an isolated copy of the article using
\code{latexmk} and \code{pdflatex}.  Optional figure regeneration
uses Matplotlib.  The standard-library runtime checks do not
require either Regina or Matplotlib.

All raw experiment reports are preserved as produced, including
their original source paths and timestamps.  The executable
drivers accept explicit paths or locate their own package root;
they do not depend on those historical absolute paths.
The complete-output digest is a deterministic correctness target.
Wall-clock times, hardware behaviour, and PDF metadata are not
expected to reproduce bit-for-bit.  The manifest distinguishes
the new proposal from unchanged source copied for standalone use.

The patch adds the sparse constructor, planar producer, independent
checker, focused tests, and research drivers beneath
\path{Topology/UnknotRecognition/fast/}.  Before applying it to a
later revision, review any changes to the native source validator,
kernel representation, integer codec, and disc-certificate API.
The package's patch was checked against its exact base.  It does
not replace the maintained dispatcher or assign a new numbered
report in the repository's editorial structure.  The TeX article
and evidence can be filed according to the repository's intake
process.

